% Generated draft: requires mathematical and source review.
\documentclass[11pt]{article}
\usepackage[margin=25mm]{geometry}
\usepackage{fontspec}
\setmainfont{Latin Modern Roman}
\usepackage{amsmath,amssymb,mathtools,unicode-math}
\setmathfont{Latin Modern Math}
\usepackage{enumitem,needspace,xurl,hyperref}
\setlength{\parindent}{0pt}
\setlength{\parskip}{5pt}
\setlength{\emergencystretch}{3em}
\newcommand{\N}{\mathbb{N}}
\newcommand{\Z}{\mathbb{Z}}
\newcommand{\Q}{\mathbb{Q}}
\newcommand{\R}{\mathbb{R}}
\newcommand{\C}{\mathbb{C}}
\newcommand{\F}{\mathbb{F}}
\newcommand{\A}{\mathbb{A}}
\newcommand{\PP}{\mathbb{P}}
\newcommand{\E}{\mathbb{E}}
\newcommand{\eps}{\varepsilon}
\newcommand{\dd}{\,\mathrm{d}}
\begin{document}
\section*{UnsolvedMath ID: 30006365}
\textbf{Source problem code:} catalogue entry
\subsection{Definitions and mathematical statement}
Let \(M\) be a smooth manifold and let \(\operatorname{Diff}^{1}(M)\) denote its \(C^{1}\) diffeomorphisms. A map \(f\in\operatorname{Diff}^{1}(M)\) is topologically transitive if, for every pair of nonempty open sets \(U,V\subseteq M\), there exists \(n\in\Z_{\geq 0}\) such that \(f^{n}(U)\cap V\ne\varnothing\). Relative to a specified \(C^{1}\) topology on \(\operatorname{Diff}^{1}(M)\), call \(f\) \(C^{1}\)-stably transitive if there is an open neighborhood \(\mathcal U\) of \(f\) such that every \(h\in\mathcal U\) is topologically transitive. The supplied excerpt also uses the term robustly transitive. A diffeomorphism has no periodic points if
\[
\forall x\in M\quad\forall n\in\Z_{\geq 1},\qquad f^{n}(x)\ne x.
\]
\textbf{Main question.} Does there exist a manifold \(M\) and a \(C^{1}\)-stably transitive diffeomorphism \(f\in\operatorname{Diff}^{1}(M)\) with no periodic points? Here stable transitivity concerns all maps in a neighborhood, whereas the absence of periodic points is required only for \(f\). The \(C^{1}\) interpretation is supported by the surrounding excerpt; the question itself does not explicitly specify a topology or differentiability class. The admissible class of manifolds, including dimension, compactness, connectedness, and boundary conditions, is not supplied. Consequently, the displayed question is a formulation template whose manifold class requires confirmation, rather than an assertion that arbitrary noncompact manifolds are intended.

\textbf{Geodesic-flow question.} For a smooth Riemannian surface \((S,g)\) with Gaussian curvature \(K_g(p)<0\) for every \(p\in S\), write
\[
T^{1}S=\{(p,v)\in TS:g_p(v,v)=1\}
\]
for its unit tangent bundle. If the unit-speed geodesic flow \((\varphi_g^{t})_{t\in\R}\) is globally defined, its time-one map is \(f_g=\varphi_g^{1}:T^{1}S\to T^{1}S\). Can such a time-one map be \(C^{1}\)-stably transitive? The source does not specify whether \(S\) is closed, complete, orientable, or without boundary, nor whether the question concerns some surface and metric or a previously fixed one. Global existence of the flow and the intended surface class must therefore be checked. This question, as supplied, does not explicitly require \(f_g\) to have no periodic points; that additional property must not be inferred solely from negative curvature.

\textbf{Unresolved surrounding formulation.} The original excerpt also asks whether an undefined map \(f\), associated with an undefined \(X\), is stably transitive, with the qualification that the answer may depend on \(X\). The supplied literature assessment describes relevant non-suspension Anosov-flow time-one maps, but the excerpt does not provide their complete hypotheses. It attributes to reference [2] the claim that a positive answer would give a robustly transitive diffeomorphism without periodic points. It also reports an affirmative answer in a suspension case, citing [9], but the referent of that answer is unclear in the extracted text. These statements cannot be combined into a precise additional conjecture from this record alone.

As historical context, the excerpt credits Bonatti--D\'iaz [1] with a 1995 construction of certain \(C^{1}\)-stably transitive perturbations of Anosov-flow time-one maps, addressing a question about diffeomorphisms isotopic to the identity. This contextual assertion is supplied attribution and is not imposed as an extra hypothesis on the main question.
\subsection{Short English statement}
The main question asks whether a diffeomorphism can remain transitive under every sufficiently small perturbation while having no periodic points itself. A related question asks whether a negatively curved surface can have a stably transitive geodesic-flow time-one map. The manifold class, surface assumptions, and parts of the surrounding Anosov-flow formulation remain unspecified. The supplied record labels the questions open; this draft does not independently verify that status.
\subsection{Sources}
\begin{itemize}
\item Local UnsolvedMath registry: permanent ID \(30006365\), exact source problem code \texttt{catalogue entry}, title \emph{Aperiodic Stable Transitivity Near Anosov Flows}, Oberwolfach Reports Open Problems collection. This draft uses the supplied complete registry record; project registry locations are \texttt{data/unsolvedmath/problems.json} and \texttt{catalog/problems.json}.
\item Supplied report citation: \emph{Dynamische Systeme (2026)}, Oberwolfach Reports, DOI \texttt{10.4171/owr/2025/31}, pp.\ 1647--1706. Supplied source URL: \url{https://doi.org/10.4171/owr/2025/31}. The citation's 2026 label differs from the 2025 report designation and proposed year; this discrepancy is unresolved.
\item Supplied literature source: \emph{Oberwolfach Report 31/2025}, \url{https://ems.press/content/serial-article-files/52240}. The exact supplied literature-check date is \texttt{2026-08-22}. The assessment that the aperiodic stable-transitivity question and the relevant non-suspension Anosov-flow time-one questions remain open is supplied dated curation, not independently verified here. The background expressly calls for independent verification of time-sensitive status and cited claims.
\item The excerpt identifies a contribution entitled \emph{Perturbation of the time-1 map of a generic volume-preserving 3-dimensional Anosov flow}, attributed to Zhiyuan Zhang, joint work with Masato Tsujii. The registry supplies no individual proposer. This contribution title alone does not establish that every listed qualification applies to each question.
\item Internal references retained from the excerpt: Bonatti--D\'iaz [1], with the supplied year 1995; [2], for the undefined \(f,X\) question and its proposed connection to aperiodicity; [6], [12, Problem 4], and [15, Problem 1], for earlier appearances of the geodesic-flow question; and [9], for the reported suspension-case remark. Full bibliographic entries are absent. A related question is attributed to Palis and Pugh, but its statement is not supplied.
\item The registry gives status \texttt{open} and statement status \texttt{corrected\_verified}, and reports that the official report was checked and notation repaired. These are supplied curation claims, not independent source verification or mathematical review performed for this draft.
\end{itemize}
\end{document}
